\documentclass[11pt]{article}

\usepackage[margin=1in]{geometry}
\usepackage{amsmath,amssymb,mathtools}
\usepackage{booktabs,tabularx,array}
\usepackage{microtype}
\usepackage{enumitem}
\usepackage{xcolor}
\usepackage{xurl}
\usepackage[hidelinks]{hyperref}

% House trust-mark palette (matches the program manuscripts).
\definecolor{kernelcolor}{RGB}{0,92,74}
\definecolor{opencolor}{RGB}{145,76,0}
\definecolor{classicalcolor}{RGB}{75,75,75}
\newcommand{\Kernel}{\textcolor{kernelcolor}{\textsc{Kernel}}}
\newcommand{\Open}{\textcolor{opencolor}{\textsc{Open}}}
\newcommand{\ImportMark}{\textcolor{classicalcolor}{\textsc{Import}}}
\newcommand{\ConjMark}{\textcolor{opencolor}{\textsc{Conjecture}}}

\newcommand{\ii}{\mathrm{i}}
\newcommand{\C}{\mathbb{C}}
\newcommand{\R}{\mathbb{R}}
\newcommand{\Lam}{\Lambda}
\newcommand{\tr}{\operatorname{tr}}
\newcommand{\Var}{\operatorname{Var}}

\setlist{itemsep=2pt,topsep=4pt}

\title{The Null-Edge Research Program\\
\large An Orientation Packet for the Non-Specialist Reader}
% RELEASE GATE: confirm the author line before circulating outside the project.
\author{Mark Schwab}
\date{July 12, 2026}

\begin{document}
\maketitle

\begin{abstract}
\noindent
This packet explains, for a reader with curiosity but no physics or
mathematics training, what the null-edge research program is trying to do,
how it works, what it has actually established so far, and what stands
between it and its goals.  The program explores a simple but radical
starting point --- that the elementary events of nature are discrete,
light-speed links called \emph{null edges}, and that mass, particle
dynamics, and perhaps the geometry of spacetime itself are bookkeeping
built on top of them.  What makes the program unusual is not the
speculation (physics has no shortage of speculative programs) but the
method: every headline mathematical claim is checked, symbol by symbol, by
a computer proof verifier, and every claim that is \emph{not} so checked
is labeled as such, with a stated test that could kill it.  This document
is deliberately honest about the boundary between what has been proved and
what is still hoped for.
\end{abstract}

\tableofcontents

\clearpage

%=====================================================================
\section{What this packet is}
%=====================================================================

The null-edge program is a research effort in mathematical physics with two
intertwined goals:

\begin{enumerate}
\item \textbf{A physics goal.}  Explore whether familiar physics --- the
  mass of particles, the Dirac equation that governs electrons, the
  symmetries of the Standard Model, and the mysterious cosmological
  constant that accelerates the universe's expansion --- can be
  \emph{reconstructed} from a strictly finite, discrete substrate whose
  only primitive ingredient is the light-speed link.
\item \textbf{A methodology goal.}  Demonstrate that speculative physics
  can be done to a higher standard of honesty than the field's norm, by
  machine-verifying every load-bearing mathematical claim and by
  pre-registering, in public, exactly which observations or theorems would
  falsify each conjecture.
\end{enumerate}

The two goals reinforce each other.  Speculative programs usually fail not
with a bang but with a slow drift into unfalsifiability; machine
verification and pre-registered kill conditions make that drift impossible
to hide.  Conversely, the discipline is only interesting if applied to
questions worth asking.

This packet has four main parts: the overarching framework and strategy
(Sections~\ref{sec:idea}--\ref{sec:framework}), the publications currently
in preparation (Section~\ref{sec:papers}), the most significant results so
far (Section~\ref{sec:results}), and the biggest remaining challenges
(Section~\ref{sec:challenges}).

%=====================================================================
\section{The idea in plain language}\label{sec:idea}
%=====================================================================

\subsection{Starting from light}

In Einstein's relativity, directions through spacetime come in three
kinds.  \emph{Timelike} directions are the paths of ordinary massive
objects: you, a planet, an electron at rest.  \emph{Spacelike} directions
connect events that cannot influence each other.  Between them sits the
knife-edge case: \emph{null} (or ``lightlike'') directions, the paths of
light itself.  A curious fact of relativity is that no time passes along a
null path.  A photon, in a precise sense, does not age.

Most approaches to fundamental physics treat massive matter as primary and
light as a special case.  The null-edge program inverts this.  Its
starting hypothesis is that the elementary event of nature is a discrete
hop at light speed --- a \emph{null edge} --- and that everything slower
than light is a composite phenomenon: light-speed motion that keeps
changing direction.

This is not a new intuition.  Feynman's ``checkerboard'' model of the
1940s pictured an electron in one space dimension as a particle zigzagging
at light speed, with the electron's mass setting the rate of direction
reversals.  The quantum-mechanical trembling motion called
\emph{zitterbewegung} points the same way.  The null-edge program asks:
what if this picture is taken seriously, made exactly finite (no
infinities, no limits taken on faith), and every step of the resulting
mathematics is proved rather than estimated?

\subsection{Mass as geometry: the area picture}

The program's central conceptual result gives the old zigzag intuition an
exact geometric form.  In the spinor language of relativity, each
light-speed motion is described by a two-component complex object $\psi$
(a \emph{null spinor}).  If a system contains several light-speed
constituents, one can ask: how far is the total from being light-speed
itself?

The answer is an \emph{area}.  For two null spinors $\psi_1$ and $\psi_2$,
the quantity $z = \psi_1\wedge\psi_2$ measures the area of the
parallelogram they span in their abstract two-dimensional complex space.
The program's flagship identity, proved and machine-checked, states that
for the combined momentum object $P$,
\[
  \det P \;=\; \sum_{i<j} \lvert \psi_i \wedge \psi_j \rvert^2 ,
\]
that is: \emph{the squared rest mass of the total is exactly the sum of
the squared areas between pairs of its light-speed constituents}.  A
bundle of light rays all pointing the same way spans no area and stays
massless.  The moment two genuinely different null directions open up, an
area appears --- and that area \emph{is} the mass.

In slogan form: \textbf{a massive particle is light going in more than one
direction, and its rest mass is the area between the directions.}  The
slogan is marketing; the determinant identity behind it is a theorem, and
the program is careful never to let the slogan claim more than the theorem
does.  In particular, the construction replaces the scalar mass input with
geometric data; it does not predict any particular mass value.

\subsection{Space, time, and counting}

The program sits in the intellectual tradition of \emph{causal set
theory}, which holds that spacetime is at bottom a discrete partial order
--- a locally finite web of ``this event can influence that one'' ---
and that geometry is recovered by the maxim \emph{order plus number equals
geometry}: the causal order supplies the shape (angles, light cones), and
simple counting of the discrete elements supplies the sizes (durations,
volumes).  A classical theorem of Malament makes the first half precise in
the continuum; the null-edge program treats the division of labor as a
standing discipline: whenever someone claims that geometry ``emerges,''
they must say which half --- the free half (shape from order) or the owed
half (scale from counting) --- their claim actually addresses.

This counting picture is also where the program touches cosmology.  If
spacetime volume is literally a count of discrete elements, then that
count fluctuates, and a celebrated heuristic (due to Sorkin and
collaborators) says the cosmological constant $\Lam$ should fluctuate with
magnitude $1/\sqrt{V}$ --- which, for the observed universe, lands
strikingly close to the tiny observed value of dark energy.  Whether that
heuristic survives exact scrutiny is one of the program's active fronts
(Sections~\ref{sec:papers} and~\ref{sec:results}).

%=====================================================================
\section{How we work: physics you can compile}\label{sec:method}
%=====================================================================

\subsection{The proof checker}

Every headline mathematical claim in this program is formalized in
\emph{Lean~4}, a proof assistant: a programming language in which one
writes not only definitions but complete logical proofs, which a small,
heavily scrutinized core program (the \emph{kernel}) then checks step by
step.  The kernel cannot be persuaded, tired out, or impressed by
reputation.  If a proof has a gap, the build fails.

Two consequences matter for a lay reader:

\begin{itemize}
\item \textbf{When we say ``proved,'' we mean machine-checked.}  A claim
  marked \Kernel{} in our manuscripts has been verified by the Lean kernel
  down to the axioms of standard mathematics.  There is no ``the details
  are left to the reader.''
\item \textbf{The computer checks the proof, not the physics.}  The kernel
  certifies that a stated theorem follows from the axioms.  It cannot
  certify that the stated theorem is the \emph{right} statement --- that
  it means what the surrounding prose says it means.  Guarding that
  boundary (``semantic alignment'') is a human and editorial
  responsibility the program takes as seriously as the proofs themselves,
  with named failure modes it audits against: vacuous hypotheses,
  trivialities dressed as depth, prose that outruns the theorem, and
  statements that are true but not the intended mathematics.
\end{itemize}

The formal development is substantial: the current publication packet
cites more than seventy headline Lean modules, each with recorded
provenance and explicitly documented conventions.

\subsection{The claim ladder}

Not everything in a physics program can be a theorem, and pretending
otherwise is its own form of dishonesty.  Program manuscripts therefore
label every claim on a fixed ladder:

\begin{description}[style=nextline]
\item[\Kernel] A theorem checked by the Lean kernel in this repository,
  with a standard, pinned axiom footprint.  The strongest currency we
  have.
\item[Kernel + evaluator] A small, explicitly disclosed class of results
  whose proof uses Lean's compiled evaluator for large finite
  computations (for instance, verifying identities of $28\times 28$
  integer matrices).  Trustworthy, but resting on slightly more machinery
  than the bare kernel; every such use is disclosed at the point of use,
  and the program routinely retrofits these down to bare-kernel proofs.
\item[\ImportMark] An established result from the literature that we rely
  on but have not re-proved.  Imports are named, cited, and checked
  against primary sources before release.
\item[\ConjMark] A pre-registered conjecture, stated with the test that
  would settle it and the condition under which we will declare it dead.
\item[Speculation] Labeled as such, always.  Interpretation and motivation
  are kept visibly separate from theorems.
\end{description}

\subsection{Guards, audits, and kill conditions}

Three further practices are worth explaining because they shape everything
in Sections~\ref{sec:papers} and~\ref{sec:results}.

\textbf{Axiom guards.}  Each flagship result is protected by a
\emph{guard file}: a tripwire compiled into the codebase that records the
exact logical assumptions of the theorem and fails the entire build if
they ever silently change --- for example, if a future edit weakened a
proof or smuggled in an extra assumption.  Honesty is enforced by the
compiler, not by memory.

\textbf{Adversarial audits.}  The program runs regular internal red-team
reviews (by independent AI agents) whose explicit job is to
find over-claims: abstract sentences without a theorem behind them,
theorems whose hypotheses no concrete example satisfies, conventions that
drifted between sources.  Findings are recorded in public scorecards, and
corrections are made in the manuscripts, not in private.

\textbf{Kill conditions.}  Conjectures are pre-registered with the result
that would falsify them, and the program's roadmap is a ladder of
\emph{gates}: pass/fail checkpoints designed so that both outcomes are
publishable.  If a hoped-for gap theorem fails, the obstruction theorem is
itself the paper.  The program has already exercised this discipline in
public: an early claim of exact Lorentz invariance for a particular finite
lattice was withdrawn when analysis showed it could not hold, and a
proposed identification of the electron's internal clock with a certain
canonical mathematical flow was not merely abandoned but \emph{disproved}
--- the refutation is now itself a recorded result.  When a claim dies
here, it dies in the open, with a death certificate.

A note on how the work is produced: proofs and drafts are developed
jointly by human direction and AI proof agents working around the clock.
This is an accelerant, not a trust basis --- nothing rests on any agent's
say-so, because every result must pass the kernel and the audit culture
described above.

%=====================================================================
\section{The framework and strategy}\label{sec:framework}
%=====================================================================

\subsection{The objects}

The program's mathematical world is deliberately austere.  Its recurring
cast:

\begin{itemize}
\item \textbf{Null edges}: discrete light-speed links; the primitive.
\item \textbf{Decorations}: finite algebraic data carried on the edges
  (phases, spinor labels, transport rules) --- the minimum extra structure
  needed to speak of interference and internal symmetry.
\item \textbf{Quantum walks}: exactly unitary, tick-by-tick update rules
  for a quantum particle hopping along null edges.  ``Exactly unitary''
  means probability is conserved to the last decimal place at every tick;
  nothing is approximate.
\item \textbf{Finite Dirac operators}: the discrete counterparts of the
  wave operator governing electrons, assembled from null hops and
  decorations, and required to reproduce the continuum Dirac physics in
  controlled limits.
\item \textbf{Counting observables}: element counts, mode counts, defect
  counts --- the quantities that, per ``order plus number,'' are supposed
  to carry all metrical information, up to and including the cosmological
  constant.
\end{itemize}

Everything is finite: finite graphs, finite matrices, finite Hilbert
spaces.  Infinities enter only when a theorem explicitly takes a limit,
and such theorems must state their convergence rates.

\subsection{The strategy: a ladder of gates}

The research plan is organized as a ladder of named gates, in roughly this
order of logical dependence:

\begin{enumerate}
\item \textbf{Kinematics} (the ``I'' gates): make the mass-as-area
  dictionary exact --- the determinant identity, its information-theoretic
  reading (mass as entanglement between null constituents; the theorem
  that a null edge has exactly zero entropy: \emph{null edges do not
  age}), and the first-order bridge to Dirac operators.  Largely landed
  and kernel-checked.
\item \textbf{Lorentz invariance} (gate L0): no fixed finite lattice can
  be exactly Lorentz-invariant, so the program treats its lattices as
  \emph{regulators} --- disposable scaffolding --- and places the ontology
  in random causal orders that are Lorentz-invariant in distribution.
  Proving the right no-go and viability theorems here is open, paper-level
  work.
\item \textbf{Operator theory and chirality} (the ``C'' gates): discrete
  space notoriously creates spurious mirror copies of particles
  (\emph{fermion doubling}) and struggles to host the handed
  (\emph{chiral}) fermions the Standard Model needs.  The program attacks
  this with the overlap/Ginsparg--Wilson technology of lattice field
  theory, kernel-checking each step, including exact bookkeeping of where
  the doublers live and what charges they carry.
\item \textbf{Interactions and dynamics} (the ``D'' and Fock-space gates):
  promote one-particle walks to many-particle fermionic dynamics with
  exact locality (strict causal cones) and exactly computed spectra; then
  ask what principle \emph{selects} the dynamics, rather than supplying it
  by hand.
\item \textbf{Gravity and entropy} (the ``Q'' gates): the conjectured
  route to gravity runs through information theory --- monotonicity
  properties of entropy along nested causal regions.  Explicitly not
  derived yet; the slogan ``gravity is data-processing'' is a gate to be
  passed, not a result to be cited.
\item \textbf{The cosmological constant} ($\Lam$ gates): make the
  ``everpresent $\Lam$'' fluctuation heuristic exact inside the framework,
  and find out whether it survives.  This front moved decisively in July
  2026 (Section~\ref{sec:results}).
\item \textbf{Algebraic matter structure} (a parallel track): formalize
  the proposed links between the octonions --- the largest of the four
  division algebras, an eight-dimensional number system in which even the
  grouping of multiplications matters --- and the Standard Model's
  symmetry group, as a verified audit trail through constructions usually
  compared only informally.
\end{enumerate}

\subsection{What we deliberately do not claim}

The program maintains a standing list of non-claims, restated in every
relevant manuscript.  As of this writing: gravity has not been derived;
no exact Lorentz invariance is claimed for any finite structure; the
observed \emph{value} of the cosmological constant is not derived (its
scale enters as an input); no absolute particle mass is predicted (the
renormalization dictionary that would connect lattice numbers to
laboratory numbers is unbuilt); and the interacting dynamics studied so
far is \emph{supplied}, not derived from a deeper principle.  Readers
should hold us to this list.

%=====================================================================
\section{The publications in preparation}\label{sec:papers}
%=====================================================================

The current packet contains five manuscripts in active preparation, plus
supporting program documents.  Table~\ref{tab:papers} gives the one-line
view; the prose below explains each.

\begin{table}[h]
\centering
\small
\begin{tabularx}{\textwidth}{@{}lXl@{}}
\toprule
Paper & Subject in one line & Status \\
\midrule
A & Rest mass derived as null-spinor area; exactly unitary Dirac walk & frozen draft, near-ready \\
C & Defects in the mass field: exact mode counting and a positional law & frozen draft \\
E & Interacting fermionic walk: exact two-particle spectrum & working draft \\
$\Lam$ & Cosmological constant: structural core, everpresent dichotomy & new draft (July 2026) \\
FB & Octonions and the Standard Model: verified algebraic audit trail & draft \\
\bottomrule
\end{tabularx}
\caption{The active manuscripts.  Additional lanes (a strict-locality
resource theorem, a changing-lattice continuum study, and a
classification-plus-selector pair) are in earlier stages.}
\label{tab:papers}
\end{table}

\subsection{Paper A: from area to Dirac gap}

\emph{Null-Spinor Area as the Rest Gap of an Exactly Unitary Dirac Walk.}
This is the flagship.  It proves the determinant identity of
Section~\ref{sec:idea}, constructs from the area datum $z$ the unique
odd Hermitian \emph{rest operator} $B_z$ with $B_z^2 = \det(P)\,\mathbf 1$,
exponentiates it into an exactly unitary quantum walk, and proves that the
walk's fixed-momentum behavior converges to genuine Dirac flow at an
explicit rate.  It further shows the \emph{phase} of the area (not just
its size) is physically visible in an exact two-site spectrum formula, and
that the same symmetry that fixes the rest operator also selects a
specific pair interaction.  Just as importantly, the paper proves exact
no-go theorems delimiting its own construction: the minimal architecture
provably retains high-frequency artifacts and cannot be repaired by the
simplest fixes.  The paper says so, with proofs.

\subsection{Paper C: the half-winding defect paper}

If the mass field $z(x)$ varies from place to place, its zeros ---
points where the two light directions momentarily align --- act as
\emph{defects}, and localized modes appear at them.  Folklore says a
topological winding number should predict these modes.  Paper C proves,
by exhaustive machine-checked classification of all sixteen sign fields
on a four-site ring, that the folklore is wrong here: two fields with
identical winding have different mode structure.  The exact law
(kernel-checked as of July 12) is \emph{positional}: what matters is
where the defect sits relative to the reflection-symmetric sites.  The
paper is thus a precision counterexample-plus-repair: familiar invariants
(determinants, naive indices, both symmetric-time windings) are proved
blind, and the correct finite criterion is proved in their place.

\subsection{Paper E: exact interacting dynamics}

Most exactly solvable models are free (non-interacting).  Paper E takes
the program's fermionic walk, adds the selected two-particle interaction
gate, and computes the composed dynamics \emph{exactly} on a four-site
ring: the two-particle spectrum reduces to named free levels plus twelve
interacting quasienergies that all solve a single rational cubic
equation.  Every step is machine-checked, with one honestly disclosed
concession: the largest matrix identities are verified by Lean's compiled
evaluator on an integer-arithmetic twin of the problem rather than by the
bare kernel.  An exact, machine-verified interacting spectrum at this
level of explicitness is a rare artifact, aimed at the quantum-simulation
and cellular-automaton communities.

\subsection{The cosmological-constant paper}

New in July 2026.  The paper formalizes the arithmetic core of the
null-edge account of $\Lam$: (i) a structural theorem that in the finite
spectral-action bookkeeping, \emph{only counting statistics} can touch the
$\Lam$ coefficient --- no deformation of the dynamics reaches it; (ii) the
everpresent scaling law, deriving the $1/\sqrt V$ fluctuation magnitude
from a Poisson counting input, routed through the framework's own edge
count; and (iii) a pre-registered dichotomy, now resolved at the
mathematical level (Section~\ref{sec:results}), between counting
statistics that keep dark-energy fluctuations alive and quieter
(``hyperuniform'') statistics that kill them.  The paper is explicit that
it does not derive the observed value, the sign, or the stochastic
dynamics of $\Lam$.

\subsection{The Furey--Baez octonion paper}

A different kind of contribution: not new physics, but a verified audit
trail through a web of algebraic constructions --- due to Furey, Baez,
Dubois-Violette, Todorov, and others --- that link the octonions to the
Standard Model's symmetry group $S(U(2)\times U(3))$.  The paper's
flagship theorem proves that the automorphisms of an explicit octonion
model fixing one chosen imaginary unit form exactly $SU(3)$, the symmetry
group of the strong force (as a precise algebraic statement, deliberately
distinguished from claims about smooth Lie groups).  Around it sit a
formalized one-generation particle package, the exact six-element kernel
of the Standard Model's gauge-group covering, charge operators with their
eigenvalue tables, and anomaly sums --- with the one conventional input
(the right-handed singlet completion) recorded as a visible,
machine-readable claim boundary rather than hidden.

\subsection{Supporting program documents}

Behind the manuscripts sit the program treatises (the NERD series) that
record the ontology, the gate ladder, failure modes, and the standing
non-claims; a living conventions file that locks signs, orientations, and
normalization choices; and a reproducibility artifact that rebuilds the
formal development and re-verifies the headline theorems from scratch.

%=====================================================================
\section{The most significant results so far}\label{sec:results}
%=====================================================================

Everything in this section marked \Kernel{} is machine-checked and
guard-pinned as described in Section~\ref{sec:method}.

\subsection{Mass is an area (\Kernel)}

The determinant identity $\det P=\sum_{i<j}|\psi_i\wedge\psi_j|^2$ holds for
any finite family of light-speed constituents.  For a \emph{pair} of edges
the rest operator obeys the scalar law $B_z^2=\det(P)\,\mathbf 1$, so the
squared rest mass of the pair is exactly the squared area between their two
spinors, and the mass term of the resulting Dirac walk is built from that
area, not inserted by hand.  For an arbitrary family the governing statement
is instead the \emph{cube law} $B^3=\mu^2 B$ on a rank-four support block;
the scalar square is the two-edge ($n=2$) specialization, not the general
law.

The same pipeline appears as the shared figure in the undergraduate and
adjacent-researcher briefs:

\begin{verbatim}
  psi_1, psi_2         two spinors, one per light-speed edge (a PAIR)
      |
      v
  P = psi_1 psi_1* + psi_2 psi_2*    (2x2 momentum matrix of the edge pair)
      |
      v
  det P = |psi_1 wedge psi_2|^2      [KERNEL]
      |
      v
  B_z^2 = det(P) * 1                 [KERNEL, pair case]
      |
      v
  mass term of the finite Dirac walk, built FROM the area
\end{verbatim}

\subsection{Null edges do not age (\Kernel)}

The information-theoretic dictionary: for the normalized two-level object
attached to a momentum, the von Neumann entropy is zero exactly when the
momentum is null.  Purity, masslessness, and ``no internal clock'' are
provably the same thing in the finite model --- turning a poetic slogan of
the framework into a theorem, and giving mass a second reading as
\emph{entanglement} between null constituents.

\subsection{An exactly solved interacting quantum automaton (\Kernel{} + disclosed evaluator)}

Paper E's headline: the interacting two-particle spectrum of the composed
fermionic walk, pinned exactly --- the characteristic polynomial factors
into named free levels and a palindromic factor whose twelve interacting
quasienergies solve one rational cubic.  Independent experts identified
this as the packet's single most broadly interesting artifact: exact,
machine-verified interacting dynamics is something very few groups in any
tradition can exhibit.

\subsection{Exact bookkeeping of fermion doubling (\Kernel)}

For the ordered $3{+}1$-dimensional walk, all zero- and $\pi$-quasienergy
crossings --- the ``doubler'' modes that discretization inevitably creates
--- are classified with their charges computed \emph{from the walk's own
symbol} at all eight nodes: charges $\pm1$, opposite at the two gaps,
summing to zero.  This is a machine-checked finite instance of the famous
Nielsen--Ninomiya obstruction, and the exact ledger one needs before any
honest attempt to evade it.

\subsection{The positional defect law (\Kernel)}

Paper C's classification: mode counts at mass-field defects on the
four-site ring are \emph{not} determined by the winding number ---
explicit counterexample pairs are proved --- but by a positional
compatibility law relative to the reflection-fixed sites, now
kernel-checked at family level, with exact multiplicities (2, 4, or 0
modes) certified for all sixteen fields.

\subsection{The dynamics selects itself, at block level (\Kernel)}

Within the static family of rest operators, the symmetry group that
preserves the construction is proved to be exactly the chiral phase
circle, and this covariance \emph{forces} the specific quartic pair
interaction used in Paper E --- resolving what had been an open selection
conjecture, at the block level.  (Whether the full time-dependent dynamics
forces the same choice remains open, and the papers say so.)

\subsection{The everpresent-$\Lam$ fork is now a theorem (\Kernel)}

The July 2026 capstone.  The cosmological-constant story hinges on how the
fundamental count $N$ fluctuates with volume: Poisson-like scatter
($\Var(N)\sim V$) sustains ``everpresent'' dark-energy fluctuations at the
observed order of magnitude, while sufficiently correlated, quiet
(\emph{hyperuniform}) counting kills them.  Skeptics could reasonably ask
whether the quiet branch was even mathematically possible for the
framework's own fermionic states.  It is: via Wick's number-variance
identity, an explicit Fermi-sea (projection) kernel is constructed with
region size $k^2$ and count variance exactly $k/4$ --- genuinely
sub-extensive \emph{and} unbounded, the non-degenerate case.  Both
branches of the dichotomy are therefore realized by explicit states, the
fork is a kernel-checked theorem, and the surviving open question is
sharply physical rather than mathematical: \emph{which count does the
dynamics make conjugate to $\Lam$ --- a thermal one (dark energy
survives) or a Fermi-sea one (it does not)?}  A companion theorem (a
general-$N$ finite uncertainty principle of Donoho--Stark type) supports
the conjugacy picture on finite registers.

\subsection{$SU(3)$ from the octonions, verified (\Kernel)}

The Furey--Baez paper's flagship: the algebraically defined automorphisms
of an explicit octonion model fixing one imaginary unit are proved
multiplicatively equivalent to $SU(3)$ --- with the equivalence upgraded
to a full group isomorphism, and the target proved equal to the standard
special-unitary group of the mathematical library.  Around it, the
one-generation package, the $\mathbb{Z}_6$ kernel computation, and the
two-sided stabilizer characterization on the exceptional Jordan algebra
give the octonion--Standard-Model literature a machine-verified
common core.

\subsection{Chiral fermions on the regulator (\Kernel, draft lane)}

On the operator-theory front, the free tetrahedral operator has a proved
spectral gap, a self-adjoint sign, an exact Ginsparg--Wilson relation, and
kernel-checked chiral (Weyl) projectors --- the statement that the
program's regulator carries handed fermions at the free level, with the
gauge-field, index, and anomaly layers (the hard part) explicitly next.

%=====================================================================
\section{The biggest remaining challenges}\label{sec:challenges}
%=====================================================================

This section is the packet's most important honesty device.  The program's
finite theorems are exact; the distance between them and established
physics is real and measured below.

\subsection{The continuum limit}

Every landed theorem lives on a finite structure.  Physics as observed is
(at least effectively) continuous.  Paper A controls fixed-momentum
convergence to Dirac flow with explicit rates.  The development now also
contains a changing momentum spacing, exhausting physical boxes, a uniform
live-walk multiplier estimate, a normalized cell-average contraction, and
an exact three-term projection bound.  The final arbitrary-$L^2$ projection
limit, its composition with the live walk, the position-space ``commuting
square'' (finite construction then limit equals continuum construction then
square), and decoupling of high-frequency artifacts remain open.  The finite
square can be exact while the limit fails; the program's own test ladder says
so and tests for it.

\subsection{Lorentz invariance}

No fixed finite lattice is exactly Lorentz-invariant, and the program has
formally withdrawn any such claim.  The intended resolution --- ontology
in Lorentz-invariant-in-distribution random causal orders, lattices
demoted to regulators --- requires two logically separate results: a
fixed-finite-support no-go under noncompact boosts, and a measure-theoretic
invariance theorem for the distribution of a random causal order.  Finite
shadows and the relevant Bombelli--Henson--Sorkin-type literature are known,
but the program has not formalized the continuum distributional theorem or
connected it to the null-edge decorations.  This is the single largest
structural debt in the framework.

\subsection{Dynamics is supplied, not derived}

The interaction gate of Paper E is \emph{selected} by symmetry within a
declared family --- a genuine theorem --- but the family itself, and the
alternation of free and interacting layers, are choices.  The static
covariance group has been classified, and its two branches now act on the
exact ordered two-channel walk (the orientation-reversing branch together
with parity).  Exhaustive classification of the discrete step and its full
$3+1$ counterpart remain open.  More deeply, a principle
(maximum-entropy equilibrium, modular-flow generation, or something else)
must select the physical state and constraints rather than merely recover a
flow after they are supplied.  That gate-D work is formulated but unproven.

\subsection{No absolute scales or continuum gravity; no physical value of $\Lam$}

Three non-claims that must eventually become claims if the program is to
be physics rather than mathematics:

\begin{itemize}
\item \textbf{Masses.}  The framework derives mass \emph{operators}, not
  mass \emph{values}.  Connecting lattice quantities to laboratory
  numbers requires a renormalization dictionary that does not yet exist.
\item \textbf{Gravity.}  Finite soldering, source, Clausius, and
  action-variation avatars are machine-checked, but the
  entropy-monotonicity route (Q gates) is still a plan rather than a
  continuum derivation.  Nothing in the program yet reproduces Einstein's
  equations.
\item \textbf{The cosmological constant.}  The fork theorem sharpened the
  question but pointedly did not answer it: identifying which count the
  dynamics conjugates to $\Lam$, the sign of the effect, and the actual
  stochastic dynamics (including confronting existing observational
  constraints on fluctuating dark energy) are the pre-registered next
  rungs.
\end{itemize}

\subsection{Chirality and anomalies with gauge fields}

The free-level chiral release is landed, but the Standard Model's handed
fermions interact with gauge fields, where lattice traditions meet their
hardest obstructions (doubling in interacting settings, index theorems,
anomaly cancellation).  The program has finite overlap/Ginsparg--Wilson,
index-integrality, inertia, and winding interfaces.  The immediate missing
bridge is an explicit gauge-coupled operator whose overlap sign and index are
proved gauge covariant and stable while the gap remains open.  Arithmetic
anomaly cancellation for a supplied Standard Model multiplet is a useful
control, not a substitute for that gauge-field theorem.  A no-go theorem may
still be what lies at the top.

\subsection{Unification of the threads}

The packet's threads rhyme --- the same cube law $X^3=(\text{area})\,X$
governs Paper A's rest operator and Paper E's interaction generator; the
octonionic structures echo the spinor constructions --- but rhyme is not
identity.  A kernel-checked corollary captures the shared cube-law shape,
and intermediate Jordan--Clifford rungs now include a Furey/exterior-Fock
correspondence and a representation-level $\mathbb Z_6$ kernel.  These
results explicitly \emph{disclaim} being a unification: no theorem yet
derives the interaction from the rest operator or the null-edge weak space
from the Jordan flag without additional choices.  The master bridge remains
a ladder of graded rungs, each with a kill condition, whose top rung is
unproven.  Claiming unification before an actual intertwiner or reconstruction
theorem is a named failure mode.

\subsection{The permanent challenge: semantic alignment}

Finally, the challenge that never closes: the kernel checks proofs, not
meanings.  Every additional theorem increases the surface on which a
correct proof of a subtly wrong statement could sit.  The audit culture
--- adversarial reviews, over-claim taxonomies, guard files, claim
labels --- is the program's standing answer, and it is a process, not a
theorem.

%=====================================================================
\section{How to read our claims, and what would change our minds}
%=====================================================================

A reader encountering the program's manuscripts should apply three rules:

\begin{enumerate}
\item \textbf{Trust the marks, not the prose.}  Sentences marked \Kernel{}
  are machine-checked theorems; unmarked prose is interpretation.  When in
  doubt, the theorem statement --- not the surrounding narrative --- is
  the claim.
\item \textbf{Check the non-claims list.}  Each manuscript states what it
  does not establish.  If a summary of our work (including an enthusiastic
  one) contradicts the non-claims list, the summary is wrong.
\item \textbf{Look for the kill conditions.}  Every conjecture worth the
  name is pre-registered with the result that would kill it.  A program
  that cannot state what would falsify it is not doing science; we invite
  readers to hold the program to its own stated tests.
\end{enumerate}

What would change our minds?  Concretely: a proof that the quiet
(hyperuniform) counting branch is dynamically forced would kill the
everpresent-$\Lam$ identification; failure of the Lorentz-ensemble gate
would confine the framework to regulator status; an obstruction at the
gauge-field chirality gate would demote the matter sector to an effective
description; and persistent failure of the continuum commuting square
would leave the program as exact finite mathematics with a physics-shaped
hole.  Each of these outcomes would be published under the same standards
as a success --- because the program's deepest bet is that in speculative
physics, \emph{verified honesty is itself the most valuable product}.

\bigskip

\noindent\textbf{Where to look next.}  The program's manuscripts (the five
papers of Section~\ref{sec:papers}), the reproducibility artifact that
rebuilds and re-verifies every headline theorem, and the public scorecards
and ledgers of the audit runs together form the packet's full evidence
base.  Nothing in this overview asks to be taken on trust; that is the
point.

\end{document}
